\section{Compatibilidad de los radios de la cola}
Es una implicación de dos estimaciones archivadas. Sea $P$ el punto de torres
de la misma máscara, $j\max_e P_e\le M_0$, brecha $r<1$ en la convención
$\|P-J_j\|_{\rm op}\le r$, y $\epsilon_A=\max_f\sum_{e\in f}P_e$.
La norma de activación es $\|H\|_\rho=\sum_U\rho^{|U|}|[U]H|$.
\emph{Media franjas}, partes 01 y 06, enuncia para $F$ la cota
\[
 \|F\|_\rho\le C\{1+j\rho^2(\epsilon_A+1/j)\},
 \qquad \rho\ge1,\quad \rho^2(\epsilon_A+1/j)\le\eta_F .
\]
La parte 01 efectúa la comparación con el escalado exacto $D$ dentro de la prueba:
$D_e/P_e=1+O(1/j)$ en el soporte. La carga en el enunciado sigue midiéndose con $P$.
\emph{Cola de activaciones}, parte 03, enuncia para $\rho\ge2$
\[
 \|\Lambda\|_\rho\le C\{j\log(ej)+\rho^2j\},
 \qquad \rho^2(\epsilon_A+1/j)\le\eta_\Lambda .
\]
Fijados los márgenes de referencia y asumidas las premisas de ambos resultados,
sirve $\eta=\min(\eta_F,\eta_\Lambda)>0$. Para cualquier $\rho\ge2$ fijo admitido
y carga dentro de ese umbral, las normas cuestan en conjunto $C_\rho j^2$.
En particular el valor absoluto total de los coeficientes de $F-\Lambda$
con $|U|>N$ es a lo sumo $C_\rho j^2\rho^{-N}$. Para
$N=\lceil4\log n/\log\rho\rceil$, $j\le n$, esto es $O_\rho(n^{-2})$.
La misma mayoración vale para una subfamilia de esos coeficientes; las
sustracciones necesarias para recuperar $V$ deben respetar sus soportes.
Esta comparación verifica los cuantificadores en el dominio común.
El peso de las llegadas fuera del dominio se controla en la prueba detenida.
Los dos enunciados sobre las normas conservan sus hipótesis, incluida la
aproximación de los menores con la cola ponderada de la mezcla en el lado exacto.
